% Chapter 6: The Measurability Layer
%
% This chapter is net-new content. The companion textbook does not cover
% the Borel-analytic separation, the WellBehavedVCMeasTarget refinement,
% the Choquet capacitability bridge, or the singleton-class witness,
% because all four were discovered during the formalization effort
% after the textbook's cutoff.
%
% The structure is driven by task-URS-manuscript measurements:
%   §6.1  Profile E (Obstruction)         -- the precise point where
%         Borel is stronger than the proof needs
%   §6.2  Profile B + C (Revelation at
%         the bridge)                     -- Choquet capacitability as
%         the classical regularity the learning-theory proof touches
%   §6.3  Profile C + D (Revelation with
%         witness-centred negative result) -- singleton class over an
%         analytic non-Borel set, full witness chain in five theorems

\chapter{The Measurability Layer}\label{chap:measurability}

Up to \Cref{ch:pac} the fundamental theorem of statistical learning
has been stated as if measurability were a side condition. It is not.
The symmetrization argument in the forward direction picks up a set on
$X^m \times X^m$ whose regularity determines whether Hoeffding's
inequality can be applied pointwise and lifted to a supremum over the
growth-function-many effective labelings. The literature, following
Krapp and Wirth~\cite{KrappWirth2024}, states this regularity as a
Borel hypothesis on the ghost-gap bad event. Every bad event in
practice is then checked against the Borel $\sigma$-algebra on the
product space, and everything goes through.

This chapter establishes that the Borel hypothesis is stronger than
the proof needs. The actual sufficient condition is weaker: the bad
event must be \emph{null-measurable} with respect to the completion
of the product measure. \Cref{sec:borel-param} identifies the point
in the symmetrization chain where the weakening applies and names the
kernel's refined hypothesis \texttt{WellBehavedVCMeasTarget}.
\Cref{sec:choquet-bridge} shows that every Borel-parameterized concept
class automatically satisfies the refinement via a classical theorem
of descriptive set theory (Choquet capacitability, Kechris
GTM~156~\cite[Theorem 30.13]{Kechris1995}). Finally,
\Cref{sec:borel-analytic-sep} exhibits a concrete concept class whose
ghost-gap bad event is analytic and null-measurable but \emph{not}
Borel, and which therefore witnesses that the refinement is strict.

None of this material appears in the companion textbook: the
discoveries are tied to the formalization work, and their home is
this chapter. The formalized chain is
\[
  \lean{singletonClassOn}\;\longrightarrow\;
  \lean{singleton_badEvent_eq_preimage_planar}\;\longrightarrow\;
  \lean{planarWitnessEvent_analytic}\;\longrightarrow\;
  \lean{planarWitnessEvent_not_measurable}\;\longrightarrow\;
  \lean{singleton_badEvent_not_measurable},
\]
five theorems in \texttt{FLT\_Proofs.Theorem.BorelAnalyticSeparation},
proved sorry-free against mathlib4 \texttt{fde0cc5}.

% ================================================================
% SECTION 6.1: BOREL PARAMETERIZATION AND THE REGULARITY GAP
% Profile E (Obstruction): diagnostic analysis of the exact point
% where the literature's Borel hypothesis is stronger than needed.
% ================================================================

\section{Borel parameterization and the symmetrization route}
\label{sec:borel-param}

The symmetrization argument at the heart of the uniform convergence
proof (Stage~2 of \Cref{sec:vc-char}) rewrites the one-sided bad event
\[
  \mathrm{Bad}_1(C, m, \eps) \;=\;
  \Bigl\{\, S \in X^m \;:\;
    \sup_{h \in C}\,
    \bigl|\, \risk_D(h) - \emprisk_S(h) \,\bigr| > \eps
  \,\Bigr\}
\]
in terms of a ghost sample $S' \sim D^m$ and the two-sided comparison
event
\[
  \mathrm{Bad}_2(C, m, \eps) \;=\;
  \Bigl\{\, (S, S') \in X^m \times X^m \;:\;
    \sup_{h \in C}\,
    \bigl|\, \emprisk_{S'}(h) - \emprisk_S(h) \,\bigr| > \eps
  \,\Bigr\}.
\]
The bound
$\Pr_S[\mathrm{Bad}_1] \leq 2\,\Pr_{S,S'}[\mathrm{Bad}_2]$
transfers control of the deviation between empirical and true error
over to the two-sample event. The union-bound-over-labelings step
that follows requires the inner event at each fixed labeling to carry
a well-defined probability: for each split of the combined $2m$
points into $S$ and $S'$, Hoeffding's inequality is applied to the
$2m$ centered Bernoulli contributions, and the resulting tail bound
is summed over the $\growth_C(2m)$ distinct effective labelings.

What does this argument actually require of $\mathrm{Bad}_2$? Three
things:

\begin{enumerate}[leftmargin=*,itemsep=3pt]
  \item \textbf{A probability for $\mathrm{Bad}_2$.} The outer
    symmetrization bound compares
    $\Pr_{S,S'}[\mathrm{Bad}_2]$ against a sum of inner tail bounds.
    For this comparison to be meaningful, $\mathrm{Bad}_2$ must be
    assigned a probability with respect to the product measure
    $D^{\otimes 2m}$.
  \item \textbf{Measurability of each slice at a fixed labeling.}
    The inner Hoeffding bound is applied to the event that a specific
    (fixed) labeling produces a large gap between the $S$-average and
    the $S'$-average. This inner event lives in the product Borel
    structure and is Borel automatically.
  \item \textbf{Closure of the probability under countable
    operations.} The supremum over $h \in C$ is rewritten as a
    finite union over the growth-function-many effective labelings on
    the combined sample, so the only closure needed is under finite
    union inside the product $\sigma$-algebra extended with the
    $D^{\otimes 2m}$-null sets.
\end{enumerate}

Requirement~(1) is exactly null-measurability with respect to the
completion of $D^{\otimes 2m}$. Requirements~(2) and~(3) are
automatic for any $\mathrm{Bad}_2$ built from a class with a Borel
parameter space. Nowhere does the proof need $\mathrm{Bad}_2$ itself
to be Borel with respect to the product Borel structure.

\begin{remark}[The Krapp and Wirth hypothesis, diagnosed]
\label{rmk:kw-diagnosis}
The natural place to plant a measurability hypothesis in the proof is
after~(1): require $\mathrm{Bad}_2$ to sit in whatever ambient
$\sigma$-algebra the argument is using. Krapp and
Wirth~\cite{KrappWirth2024} plant it inside the product Borel
$\sigma$-algebra. That is a strictly stronger requirement than~(1)
alone. The gap between ``Borel on the product space'' and
``null-measurable for every product measure built from $D$'' is
invisible in every Borel-parameterized example, where both happen to
hold, but it is a real gap in general.
\end{remark}

The kernel names the refined hypothesis.

\begin{defn}[WellBehavedVCMeasTarget]
\label{def:wbvc-meas-target}
  \lean{WellBehavedVCMeasTarget}
  \leanok
A concept class $C$ over a measurable space $(X, \Sigma)$ satisfies
\emph{WellBehavedVCMeasTarget} if, for every sample size $m$, every
gap threshold $\eps$, and every probability distribution $D$ on $X$,
the two-sided ghost-gap bad event $\mathrm{Bad}_2(C, m, \eps)$ is a
null-measurable subset of
$(X^m \times X^m,\, \Sigma^{\otimes 2m})$ with respect to the product
measure $D^{\otimes 2m}$. Here null-measurability means that the
event differs from a set in $\Sigma^{\otimes 2m}$ by a
$D^{\otimes 2m}$-null set, equivalently, it sits in the completion of
$\Sigma^{\otimes 2m}$ under $D^{\otimes 2m}$.
\end{defn}

\begin{defn}[KrappWirthWellBehaved]
\label{def:krapp-wirth}
  \lean{KrappWirthWellBehaved}
  \leanok
A concept class $C$ satisfies the \emph{KrappWirthWellBehaved}
hypothesis of~\cite{KrappWirth2024} if, for every $m$ and $\eps$, the
two-sided ghost-gap bad event $\mathrm{Bad}_2(C, m, \eps)$ is a Borel
subset of $X^m \times X^m$ in the product Borel structure.
\end{defn}

\begin{prop}[KrappWirth is stronger than WellBehavedVCMeasTarget]
\label{prop:kw-implies-wbvc}
  \lean{KrappWirthWellBehaved.toWellBehavedVC}
  \leanok
  \uses{def:wbvc-meas-target,def:krapp-wirth}
Every Borel set is null-measurable with respect to any completion of
a Borel probability measure. Consequently any concept class carrying
\texttt{KrappWirthWellBehaved} also carries
\texttt{WellBehavedVCMeasTarget}.
\end{prop}

\begin{proof}
$\mathcal{B} \subseteq \mathcal{B}_\mu^\ast$ for every Borel
probability measure $\mu$, because the completion only enlarges the
$\sigma$-algebra by $\mu$-null sets. Apply this to the product
structure $\Sigma^{\otimes 2m}$ with the product measure
$D^{\otimes 2m}$ and the statement follows immediately.
\end{proof}

What \Cref{prop:kw-implies-wbvc} leaves open is the reverse. The
reverse direction is the subject of \Cref{sec:borel-analytic-sep}. We
prepare the ground by introducing the regularity the kernel uses in
practice to produce instances of \texttt{WellBehavedVCMeasTarget}.

\begin{defn}[Borel-parameterized concept class]
\label{def:borel-param}
A concept class $C \subseteq \{0,1\}^X$ over a measurable space
$(X, \Sigma)$ is \emph{Borel-parameterized} by a standard Borel space
$(\Theta, \mathcal{B}_\Theta)$ if there exists a jointly measurable
evaluation map
$\mathrm{eval} : \Theta \times X \to \{0,1\}$
such that
$C = \bigl\{\,\mathrm{eval}(\theta, \cdot)\,:\,\theta \in \Theta\,\bigr\}$.
\end{defn}

Borel parameterization is the standard regularity assumption in
descriptive-set-theory approaches to learning theory. Every concept
class considered in practice admits a Borel parameterization: finite
classes, neural networks with measurable weights, decision trees of
bounded depth, thresholds on $\R$, halfspaces, axis-aligned rectangles,
and so on. When the parameter space $\Theta$ is discrete or finite,
joint measurability is automatic. When $\Theta = \R^k$ for some
$k$, joint measurability amounts to continuity of $\mathrm{eval}$ in
$\theta$ for each fixed $x$ together with measurability of
$\mathrm{eval}$ in $x$ for each fixed $\theta$. The kernel's
\texttt{WellBehavedVCMeasTarget} setting captures exactly this
regularity, and the next section proves that the refinement is
automatic for every such class.

\begin{defn}[Parameterized ghost-gap bad event]
\label{def:param-bad-event}
  \lean{paramBadEvent}
  \leanok
  \uses{def:borel-param}
For a Borel-parameterized class with evaluation map $\mathrm{eval}$,
sample size $m$, and gap threshold $\eps$, the \emph{parameterized
ghost-gap bad event} is
\[
  \mathrm{Bad}_{\mathrm{eval}}(m, \eps)
  \;=\;
  \bigcup_{\theta \in \Theta}
  \Bigl\{\, (S, S') \in X^m \times X^m \;:\;
    \bigl| \emprisk_{S'}(\mathrm{eval}(\theta,\cdot))
      - \emprisk_S(\mathrm{eval}(\theta,\cdot)) \bigr| > \eps
  \,\Bigr\}.
\]
Equivalently, $\mathrm{Bad}_{\mathrm{eval}}(m, \eps)$ is the
projection onto the last two factors of a jointly Borel subset of
$\Theta \times X^m \times X^m$.
\end{defn}

The ``projection of a jointly Borel event'' reading is the key. A
projection of a Borel set onto a product factor is not Borel in
general. It is, however, always analytic. This observation is what
connects the kernel's refinement to classical descriptive set theory,
and it is the topic of \Cref{sec:choquet-bridge}.

% ================================================================
% SECTION 6.2: THE ANALYTIC MEASURABILITY BRIDGE
% Profile B + C: load-bearing bridge theorem with a revelation turn
% at the point where Choquet capacitability, a pure descriptive set
% theory result, exactly matches the regularity the learning theory
% argument needs.
% ================================================================

\section{The analytic measurability bridge}
\label{sec:choquet-bridge}

Analytic sets are the sets obtainable as continuous images of Polish
spaces, equivalently, the projections of Borel sets in a product of
Polish spaces, equivalently, the result of the Souslin operation
applied to a Borel scheme. Every Borel set is analytic, the converse
fails in the strongest possible way (there exist analytic sets of
every cardinality up to the continuum that are not Borel), and the
class of analytic sets is strictly larger than the Borel class while
staying manageable under projection. This is classical, and
Kechris~\cite[Chapter~14]{Kechris1995} is the standard modern
reference.

\begin{lemma}[Projections of Borel sets are analytic]
\label{lem:borel-proj-analytic}
Let $Y$ and $Z$ be standard Borel spaces and let $B \subseteq Y \times Z$
be Borel. The projection $\pi_Z(B) = \{\,z : \exists y,\ (y,z) \in B\,\}$
is an analytic subset of $Z$. This is the defining property of
analytic sets in the Kechris formulation.
\end{lemma}

\begin{cor}[The parameterized bad event is analytic]
\label{cor:param-bad-analytic}
  \uses{def:borel-param,def:param-bad-event,lem:borel-proj-analytic}
For every Borel-parameterized concept class $C$ with parameter space
$\Theta$, every $m$, and every $\eps > 0$, the parameterized ghost-gap
bad event $\mathrm{Bad}_{\mathrm{eval}}(m, \eps)$ is an analytic
subset of $X^m \times X^m$ whenever $X$ is standard Borel.
\end{cor}

\begin{proof}
The condition
$|\emprisk_{S'}(\mathrm{eval}(\theta,\cdot))
  - \emprisk_S(\mathrm{eval}(\theta,\cdot))| > \eps$
is jointly Borel in $(\theta, S, S') \in \Theta \times X^m \times X^m$
because $\mathrm{eval}$ is jointly measurable and the empirical risks
are finite sums of indicator evaluations. The bad event
$\mathrm{Bad}_{\mathrm{eval}}(m, \eps)$ is the projection of this
Borel set onto the last two factors, hence analytic by
\Cref{lem:borel-proj-analytic}.
\end{proof}

\Cref{cor:param-bad-analytic} reduces the measurability question for
the bad event to a single question: \emph{are analytic sets
null-measurable with respect to every Borel probability measure?}
That question was answered in 1955 by Gustave Choquet. The answer is
yes, and the proof route is the capacitability theorem.

\begin{thm}[Choquet capacitability, {\cite[Theorem 30.13]{Kechris1995}}]
\label{thm:choquet-capacity}
  \lean{MeasureTheory.AnalyticSet.compactCap_eq}
  \leanok
Let $X$ be a Polish space and let $\mu$ be a Borel probability
measure on $X$. For every analytic set $A \subseteq X$,
\[
  \mu^*(A) \;=\; \sup\,\bigl\{\, \mu(K)\,:\,K \subseteq A,\ K\text{ compact}\,\bigr\},
\]
where $\mu^*$ denotes the outer measure induced by $\mu$. In
particular, $A$ is $\mu$-capacitable: its outer measure equals its
inner measure, and $A$ lies in the completion of the Borel
$\sigma$-algebra with respect to $\mu$.
\end{thm}

\begin{proof}[Sketch]
The argument proceeds by showing that $\mu^*$ extended to analytic
sets is a Choquet capacity: it is countably subadditive, monotone,
and continuous from below on increasing unions, and continuous from
above on decreasing intersections of compact sets. Any Choquet
capacity is capacitable on analytic sets by an alternating game
between approximating unions and intersections of compact sets. The
compact approximations then converge in outer measure to $\mu^*(A)$
from below, and the Borel completion absorbs the gap.
The full formalization runs 416 lines in
\texttt{FLT\_Proofs.PureMath.ChoquetCapacity} and is independent of
learning theory; it is written in a form suitable for contribution to
Mathlib as a standalone result. The blueprint treats it as a black
box from here on.
\end{proof}

\begin{cor}[Analytic sets are null-measurable]
\label{cor:analytic-null-measurable}
  \lean{analyticSet_nullMeasurableSet}
  \leanok
  \uses{thm:choquet-capacity}
Every analytic set in a Polish space is null-measurable with respect
to every Borel probability measure.
\end{cor}

\begin{proof}
By \Cref{thm:choquet-capacity}, the outer and inner measures of an
analytic set agree, which is equivalent to the set lying in the
$\mu$-completion of the Borel $\sigma$-algebra.
\end{proof}

The bridge closes with a one-line composition:
\Cref{cor:param-bad-analytic} says the bad event is analytic,
\Cref{cor:analytic-null-measurable} says analytic sets are
null-measurable, and null-measurability of the bad event is exactly
\texttt{WellBehavedVCMeasTarget}.

\begin{prop}[Borel-parameterized implies WellBehavedVCMeasTarget]
\label{prop:borel-param-implies-wbvc}
  \lean{analyticSet_nullMeasurableSet_ghostPairs}
  \leanok
  \uses{def:borel-param,cor:param-bad-analytic,cor:analytic-null-measurable,def:wbvc-meas-target}
Every Borel-parameterized concept class over a standard Borel space
$X$ satisfies \texttt{WellBehavedVCMeasTarget}.
\end{prop}

The fundamental theorem of statistical learning (\Cref{thm:vc-char})
therefore applies, in the kernel's formulation, to every
Borel-parameterized concept class, including the classes for which
Krapp and Wirth originally stated it under the stronger Borel
hypothesis. The refinement is automatic for practical examples. The
question that remains is whether it is \emph{strictly} weaker, and it
is to that question we now turn.

% ================================================================
% SECTION 6.3: THE BOREL-ANALYTIC SEPARATION THEOREM
% Profile C + D: Revelation (the gap exists) with witness-centered
% negative-result construction (singleton class over an analytic
% non-Borel set realizes the gap). Three expository moves:
%   1. Historical box on Krapp-Wirth 2024 framing
%   2. Singleton class construction with TikZ planar witness diagram
%   3. Separation environment wrapping the main theorem
% ================================================================

\section{The Borel-analytic separation theorem}
\label{sec:borel-analytic-sep}

\Cref{prop:borel-param-implies-wbvc} says that every reasonable
concept class satisfies the kernel's measurability hypothesis via a
classical descriptive-set-theory route.  What the bridge does not
rule out is that the underlying refinement might be cosmetic: perhaps
every class that admits any kind of measurability certificate also
happens to admit a Borel one, and the \texttt{WellBehavedVCMeasTarget}
$\subsetneq$ \texttt{KrappWirthWellBehaved} inclusion is vacuous in
content.  This section rules that out with a witness.

\begin{historical}
\textbf{Krapp and Wirth 2024.}  Shortly before this formalization
effort began, Lothar Sebastian Krapp and Laura Wirth posted
\emph{Measurability in the Fundamental Theorem of Statistical
Learning} (arXiv:2410.10243, October 2024).  The paper performs
exactly the audit this chapter is built around: it scrutinizes the
standard proofs of the fundamental theorem of statistical learning in
the agnostic setting, identifies the measurability assumptions that
are tacitly imposed, and extracts a self-contained proof that makes
those assumptions precise.  The paper's measurability hypothesis on
the ghost-gap bad event is that it be \emph{Borel} in the product
Borel structure.  The refinement in this chapter observes that the
Borel hypothesis is strictly stronger than the proof needs, and this
section constructs a concept class that realizes the strict gap.  The
Krapp-Wirth paper is correct: everything stated in it under the
Borel hypothesis also holds under the weaker null-measurability
hypothesis used here.  What the present chapter adds is that the
Borel hypothesis is a genuine commitment, not a free choice, and the
commitment can be relaxed.
\end{historical}

The witness is a family of very thin concept classes on the real line.
For each set $A \subseteq \R$, let $\mathbf{1}_{\{a\}}$ denote the
indicator function of the singleton $\{a\}$, regarded as an element of
$\{0,1\}^\R$.

\begin{defn}[Singleton class over a set]
\label{def:singleton-class}
  \lean{singletonClassOn}
  \leanok
For $A \subseteq \R$, define the \emph{singleton class over $A$} as
\[
  C_A \;=\; \{\,\mathbf{0}\,\}\,\cup\,
            \bigl\{\,\mathbf{1}_{\{a\}}\,:\,a \in A\,\bigr\},
\]
where $\mathbf{0}$ is the everywhere-false concept
($\mathbf{0}(x) = 0$ for every $x \in \R$).
\end{defn}

A singleton class looks unremarkable from the outside. It has VC
dimension at most $1$, which makes it PAC learnable under any
reasonable measurability hypothesis (one can learn the identity of
the non-trivial hypothesis from a single labeled example). Its
individual hypotheses are all individually measurable in the strongest
possible sense: the zero concept is constant, and each
$\mathbf{1}_{\{a\}}$ is the characteristic function of a Borel
singleton.

\begin{lemma}[Every hypothesis in $C_A$ is measurable]
\label{lem:singleton-class-measurable}
  \lean{singletonClassOn_measurable}
  \leanok
  \uses{def:singleton-class}
For every $A \subseteq \R$, every hypothesis $h \in C_A$ is Borel
measurable.
\end{lemma}

The surprise is that individual measurability of every hypothesis
does \emph{not} lift to measurability of the ghost-gap bad event
for the class as a whole. The bad event is an indexed union over
$C_A$, and the index set $A$ can be arbitrarily irregular. The
kernel formalizes the reduction from the bad event to a planar
witness.

\begin{defn}[Planar witness event]
\label{def:planar-witness}
  \lean{planarWitnessEvent}
  \leanok
For $A \subseteq \R$, the \emph{planar witness event} is the subset
of $\R^2$ given by
\[
  W_A \;=\; \bigl\{\,(x, y) \in \R^2 \;:\; y \in A\ \text{and}\ x \neq y\,\bigr\}.
\]
\end{defn}

\begin{figure}[ht]
\centering
\begin{tikzpicture}[scale=1.0,
  axis/.style={-{Stealth[length=2.5mm]}, thick, black},
  diag/.style={dashed, thick, notegray},
  fiber/.style={thick, thmred},
  setA/.style={thick, defblue, fill=defblue!10},
  lbl/.style={font=\small}
]
  % Axes
  \draw[axis] (-0.3, 0) -- (5.2, 0) node[lbl, right] {$x$};
  \draw[axis] (0, -0.3) -- (0, 4.5) node[lbl, above] {$y$};

  % Diagonal y = x
  \draw[diag] (-0.2, -0.2) -- (4.6, 4.6) node[lbl, right, pos=0.9] {$y = x$};

  % A as a shaded horizontal strip at y = a1, a2, a3
  \foreach \ay in {1.2, 2.3, 3.6} {
    \draw[fiber] (-0.05, \ay) -- (5.05, \ay);
  }
  \node[lbl, text=thmred, anchor=west] at (5.1, 1.2) {$y = a_1$};
  \node[lbl, text=thmred, anchor=west] at (5.1, 2.3) {$y = a_2$};
  \node[lbl, text=thmred, anchor=west] at (5.1, 3.6) {$y = a_3$};

  % Mark the points removed from each fiber: x = y = a_i
  \foreach \ay in {1.2, 2.3, 3.6} {
    \fill[white, draw=thmred, thick] (\ay, \ay) circle (2.5pt);
  }

  % Planar witness label
  \node[lbl, defblue, anchor=south west] at (0.2, 3.9)
    {$W_A = \{(x,y)\,:\,y \in A,\ x \neq y\}$};
  \node[lbl, notegray, anchor=south west] at (0.2, 3.5)
    {(three sample fibers shown)};
\end{tikzpicture}
\caption{The planar witness $W_A$ depicted for three values $a_1, a_2,
a_3 \in A$.  Each fiber $\{y = a_i\}$ contributes the horizontal line
at height $a_i$, punctured at the diagonal point $(a_i, a_i)$.  The
full set $W_A$ is the union of these punctured horizontal lines over
all $y \in A$.  When $A$ is analytic, $W_A$ is analytic. When $A$ is
not Borel, $W_A$ is not Borel either.}
\label{fig:planar-witness}
\end{figure}

\begin{lemma}[Bad event reduces to planar-witness preimage]
\label{lem:bad-eq-preimage}
  \lean{singleton_badEvent_eq_preimage_planar}
  \leanok
  \uses{def:singleton-class,def:planar-witness}
For every $A \subseteq \R$, the one-step ghost-gap bad event
$\mathrm{Bad}(C_A, 1, 1/2)$ on $(\mathrm{Fin}\,1 \to \R) \times
(\mathrm{Fin}\,1 \to \R)$ equals the preimage of the planar witness
$W_A$ under the sample-pair projection
$(s, s') \mapsto (s(0), s'(0))$.
\end{lemma}

The reduction replaces a class-indexed existential quantifier with a
fixed planar set, trading one measurability problem for another. The
remaining problem is to understand the measurability of the planar
set.

\begin{thm}[$W_A$ is analytic for analytic $A$]
\label{thm:planar-analytic}
  \lean{planarWitnessEvent_analytic}
  \leanok
  \uses{def:planar-witness}
If $A \subseteq \R$ is analytic then $W_A \subseteq \R^2$ is analytic.
\end{thm}

\begin{proof}
The set $W_A$ is the intersection of the Borel set
$\{(x, y) : x \neq y\}$ with the cylinder $\R \times A$. Analytic
sets are closed under intersection with Borel sets, so $W_A$ is
analytic.
\end{proof}

\begin{thm}[$W_A$ is not Borel for some analytic non-Borel $A$]
\label{thm:planar-not-borel}
  \lean{planarWitnessEvent_not_measurable}
  \leanok
  \uses{def:planar-witness}
There exists an analytic set $A \subseteq \R$ such that $W_A$ is
analytic but \emph{not} Borel. The classical construction takes any
complete analytic subset of $\R$, for example the set of codes of
ill-founded trees on $\N$ under a standard encoding; this set is
analytic and not Borel. The projection
$A \mapsto W_A$ preserves the non-Borel property in the following
sense: if $W_A$ were Borel, then the fiber of $W_A$ over any
$x_0 \notin A$ would be Borel, but that fiber is exactly $A$, a
contradiction.
\end{thm}

With $C_A$ defined, the planar reduction established, and the planar
witness classified by descriptive set theory, the main separation
theorem assembles in a single line of logic.

\begin{separation}
\textbf{Borel-analytic separation
(\Cref{cor:borel-analytic-separation}).}

\medskip\noindent
\emph{Statement.} There exists a concept class $C \subseteq \{0,1\}^\R$
such that $C$ carries \texttt{WellBehavedVCMeasTarget} but $C$ does
\emph{not} carry \texttt{KrappWirthWellBehaved}.

\medskip\noindent
\emph{Witness.}
Let $A \subseteq \R$ be any analytic non-Borel set, for example the
set of codes of ill-founded trees as in
\Cref{thm:planar-not-borel}. Take $C = C_A$, the singleton class over
$A$.

\medskip\noindent
\emph{Why the witness works.}  Every hypothesis in $C_A$ is
individually Borel measurable
(\Cref{lem:singleton-class-measurable}), so the class has the
individual regularity one might expect.  The ghost-gap bad event
coincides, by
\Cref{lem:bad-eq-preimage}, with the preimage of $W_A$ under the
sample-pair projection.  The planar witness $W_A$ is analytic by
\Cref{thm:planar-analytic} and is not Borel by
\Cref{thm:planar-not-borel}.  The sample-pair projection is a
homeomorphism onto its image, so the bad event is itself analytic and
not Borel.  Analytic sets are null-measurable by
\Cref{cor:analytic-null-measurable}, giving
\texttt{WellBehavedVCMeasTarget}.  Borel sets are not: the bad event
fails \texttt{KrappWirthWellBehaved} by the non-Borel witness.

\medskip\noindent
\emph{What the witness exploits.}  Individual measurability of the
hypotheses in a class does not lift to measurability of the
class-indexed bad event, because the bad event is an uncountable union
whose indexing set can be as wild as the analytic hierarchy permits.
The Borel $\sigma$-algebra is stable under countable unions but not
under analytic-indexed unions; the completion under any Borel measure
is exactly the extra layer needed to accommodate such unions.  The
singleton class is the minimal construction that exposes this gap: it
has VC dimension~$1$, a single non-trivial hypothesis per index, and a
planar bad event whose irregularity is inherited entirely from $A$.

\hfill\cite{KrappWirth2024},\,\cite{Kechris1995}
\end{separation}

\begin{cor}[Borel-analytic separation]
\label{cor:borel-analytic-separation}
  \lean{singleton_badEvent_not_measurable}
  \leanok
  \uses{lem:bad-eq-preimage,thm:planar-analytic,thm:planar-not-borel,cor:analytic-null-measurable,def:wbvc-meas-target,def:krapp-wirth}
There exists a concept class $C_A \subseteq \{0,1\}^\R$ such that
$C_A$ satisfies \texttt{WellBehavedVCMeasTarget} but fails
\texttt{KrappWirthWellBehaved}. Consequently the fundamental theorem
of statistical learning (\Cref{thm:vc-char}) holds for $C_A$ under
the kernel's NullMeasurableSet refinement and does not hold for
$C_A$ under the Krapp-Wirth Borel hypothesis.
\end{cor}

\begin{remark}[What the refinement changes and what it does not]
The refinement changes the shape of the measurability hypothesis
used in the fundamental theorem from a set-theoretic constraint on
the class (Borel bad event) to a measure-theoretic constraint
(null-measurable bad event). For every class that admits a Borel
parameterization, both hypotheses hold and the two versions of the
theorem are interchangeable. For classes outside Borel
parameterization, only the null-measurability version applies. The
Krapp-Wirth paper and this kernel therefore prove the same theorem
on the same universe of practically interesting classes; the
difference is at the edge of the universe, where concept classes
over analytic non-Borel sets live. That edge is sparse in practice
but non-empty, and
\Cref{cor:borel-analytic-separation} exhibits it in full.
\end{remark}

The chain of five theorems formalized in
\texttt{FLT\_Proofs.Theorem.BorelAnalyticSeparation} thus completes
the chapter. Every step is sorry-free, cross-linked with
\texttt{FLT\_Proofs.Complexity.BorelAnalyticBridge} on the analytic
measurability side, and cited from
\texttt{FLT\_Proofs.Complexity.Measurability} where the
\texttt{WellBehavedVCMeasTarget} typeclass lives.
